% =====================================================================
\chapter{Los números reales}
\label{tema:1}
% =====================================================================

Este tema fija las reglas de juego de todo el módulo: qué se puede hacer con los números reales y por qué. En lugar de dar por sabidas las propiedades de la aritmética, se parte de una lista corta de axiomas, la de un \emph{cuerpo ordenado}, y se deducen de ella las reglas habituales. La presentación sigue los capítulos~1 y~2 de Spivak~\cite{spivak}; las secciones de potencias y radicales siguen los §1.1--1.3 de OpenStax~\cite{openstax}.

\section{Conjuntos numéricos}

Se trabaja con la cadena de inclusiones $\N\subset\Z\subset\Q\subset\R$, donde $\N=\{1,2,3,\dots\}$ son los naturales, $\Z$ los enteros, $\Q=\{p/q: p,q\in\Z,\ q\neq 0\}$ los racionales y $\R$ los reales. Un real que no es racional se llama \emph{irracional}.

\begin{proposicion}\label{prop:decimal-periodico}
Un número real es racional si y solo si su expresión decimal es finita o periódica.
\end{proposicion}

\addtocounter{definicion}{1}%

\section{Axiomas de cuerpo}

\begin{definicion}[Cuerpo]\label{def:cuerpo}
Un \emph{cuerpo} es un conjunto $K$ con dos operaciones, suma $(a,b)\mapsto a+b$ y producto $(a,b)\mapsto ab$, que cumplen, para todos $a,b,c\in K$:
\begin{enumerate}[label=(C\arabic*)]
  \item asociatividad: $(a+b)+c=a+(b+c)$ y $(ab)c=a(bc)$;
  \item conmutatividad: $a+b=b+a$ y $ab=ba$;
  \item elementos neutros: existen $0,1\in K$, con $0\neq 1$, tales que $a+0=a$ y $a\cdot 1=a$;
  \item opuesto: para cada $a$ existe $-a$ con $a+(-a)=0$;
  \item inverso: para cada $a\neq 0$ existe $a^{-1}$ con $aa^{-1}=1$;
  \item distributividad: $a(b+c)=ab+ac$.
\end{enumerate}
Se definen la resta $a-b=a+(-b)$ y, si $b\neq 0$, la división $a/b=ab^{-1}$.
\end{definicion}

$\Q$ y $\R$ son cuerpos; $\Z$ no lo es, porque $2$ no tiene inverso en $\Z$. Todas las reglas de la aritmética escolar se deducen de (C1)--(C6).

\begin{proposicion}[Consecuencias de los axiomas de cuerpo]\label{prop:cuerpo}
Para todos $a,b,x$ de un cuerpo:
\begin{enumerate}
  \item si $a+x=a$, entonces $x=0$ (unicidad del neutro); si $a+x=0$, entonces $x=-a$ (unicidad del opuesto);
  \item $a\cdot 0=0$;
  \item $(-a)b=-(ab)$ y $(-a)(-b)=ab$;
  \item si $ab=0$, entonces $a=0$ o $b=0$.
\end{enumerate}
\end{proposicion}

El apartado~(iv) es el que justifica resolver $(x-2)(x+3)=0$ igualando a cero cada factor.

\section{Axiomas de orden}

\begin{definicion}[Cuerpo ordenado]\label{def:orden}
Un cuerpo $K$ es \emph{ordenado} si existe un subconjunto $P\subset K$, el de los \emph{positivos}, tal que:
\begin{enumerate}[label=(O\arabic*)]
  \item tricotomía: para cada $a\in K$ se cumple exactamente una de $a=0$, $a\in P$, $-a\in P$;
  \item si $a,b\in P$, entonces $a+b\in P$;
  \item si $a,b\in P$, entonces $ab\in P$.
\end{enumerate}
Se escribe $a<b$ (o $b>a$) si $b-a\in P$, y $a\le b$ si $a<b$ o $a=b$.
\end{definicion}

\begin{proposicion}[Reglas de las desigualdades]\label{prop:orden}
En un cuerpo ordenado, para todos $a,b,c$:
\begin{enumerate}
  \item si $a<b$ y $b<c$, entonces $a<c$;
  \item si $a<b$, entonces $a+c<b+c$;
  \item si $a<b$ y $c>0$, entonces $ac<bc$;
  \item si $a\neq 0$, entonces $a^2>0$; en particular, $1>0$.
\end{enumerate}
\end{proposicion}

\begin{observacion}[La completitud]
$\Q$ y $\R$ son ambos cuerpos ordenados; lo que los distingue es el \emph{axioma del supremo}: todo subconjunto no vacío de $\R$ acotado superiormente tiene una cota superior mínima~\cite[cap.~8]{spivak}. Gracias a él existen en $\R$ números como $\sqrt 2$, que no están en $\Q$ (\cref{teo:raiz2}). En este módulo el axioma solo se usa implícitamente, al admitir que existen las raíces $n$-ésimas y los logaritmos de los reales positivos.
\end{observacion}

\section{Números irracionales}

\begin{lema}\label{lem:par}
Sea $n\in\Z$. Si $n^2$ es par, entonces $n$ es par.
\end{lema}

\begin{teorema}\label{teo:raiz2}
$\sqrt2$ es irracional.
\end{teorema}

\section{Potencias, radicales e intervalos}

Para $a\in\R$ y $n\in\N$ se define $a^n=a\cdots a$ ($n$ factores), $a^0=1$ si $a\neq 0$, y $a^{-n}=1/a^n$. Para $a>0$ se define $a^{1/n}=\sqrt[n]{a}$ como el único real positivo cuya potencia $n$-ésima es $a$, y $a^{m/n}=(a^{1/n})^m$.

\begin{proposicion}[Leyes de los exponentes]\label{prop:exponentes}
Para $a,b>0$ y $r,s\in\Q$: $a^{r}a^{s}=a^{r+s}$, $(a^{r})^{s}=a^{rs}$, $(ab)^{r}=a^{r}b^{r}$ y $a^{-r}=1/a^{r}$. Para exponentes enteros, las tres primeras valen también con $a,b\neq 0$.
\end{proposicion}

Para exponentes naturales, las leyes se prueban por inducción sobre el exponente; para exponentes racionales se reducen a ese caso usando la unicidad de la raíz $n$-ésima~\cite[§1.2--1.3]{openstax}.

\begin{definicion}[Intervalos]
Para $a<b$: $[a,b]=\{x: a\le x\le b\}$, $(a,b)=\{x: a<x<b\}$, y análogamente $[a,b)$, $(a,b]$, $(a,\infty)$, $(-\infty,b]$, etc. Los símbolos $\pm\infty$ no son números reales y nunca llevan corchete cerrado.
\end{definicion}

\addtocounter{definicion}{1}%

\addtocounter{definicion}{1}%

